% Laget automatisk av elevark/lag_elevark.py fra programmering/python-horisontalt-kast.html. Endringer her blir overskrevet.
\documentclass[11pt]{article}
\usepackage{elevark}
\setlength{\svarhoyde}{1.4cm}
\begin{document}
\elevarkhode{Fysikk 2 · Bevegelse og krefter i to dimensjoner}{Horisontalt kast med Eulers metode}{https://fysikklab.no/programmering/python-horisontalt-kast.html}{Koden på nett}

En ball triller ut over kanten av et bord som er 0,80~m høyt, med farten 1,3~m/s. Hvor langt fra bordet lander den? Her regner vi det ut med små tidssteg, og bruker numpy-vektorer for å regne med $x$ og $y$ samtidig.

\section*{Oppgaver}

\begin{oppgaveliste}
\oppgave \fase{test}{Kjør}Kjør steg 1 og 2. Hvor godt stemmer simuleringen med de eksakte svarene?
\svar
\oppgave \fase{forklar}{Forklar}Hvorfor endrer ikke x-komponenten av farten seg? Hva betyr det at fartsvektoren er [1.3, −3.924]?
\svar
\oppgave \fase{test}{Undersøk}Prøv $dt = 0{,}05$, 0,02 og 0,005. Hvordan påvirker tidssteget svaret? Hvorfor stopper løkka alltid litt for sent?
\svar
\oppgave \fase{test}{Endre}Gjør kastet skrått: Ballen kastes fra bordkanten med 4,0~m/s, 45° over horisontalen. Hvordan må startfarten skrives som en vektor? Hvor langt fra bordet lander ballen?
\svar
\oppgave \fase{utfordring}{Utfordring}Legg til luftmotstand: $a = g - (k/m)|v| \cdot v$, der $|v|$ er farten (np.linalg.norm). Bruk $k/m = 0{,}05\,\mathrm{m^{-1}}$. Hvordan endres landingspunktet? Hvorfor er det mye enklere å gjøre dette numerisk enn for hånd?
\svar
\end{oppgaveliste}

\end{document}
